\leanverified{paper\_conclusion\_kernel\_rh\_q4\_minimal\_threshold}
\begin{theorem}[\texorpdfstring{$q=4$}{q=4} 是二次观测失去平方根稳定性的最小越线阶]\label{thm:conclusion-kernel-rh-q4-minimal-threshold}
沿用注 \ref{rem:pom-collision-kernel-rh-phase} 与定理 \ref{thm:pom-deltaq-mean-square-rh-criterion} 的记号，记
\[
\delta_q:=\log\!\left(\frac{\Lambda_q^2}{r_q}\right),
\qquad
E_q(n):=\Tr(A_q^n)-r_q^n,
\qquad
\mathcal S_q:=\sum_{n\ge 1}\frac{|E_q(n)|^2}{r_q^n}.
\]
则在已核验范围 \(2\le q\le 23\) 内存在 sharp 相变
\[
\delta_2<0,\qquad \delta_3<0,\qquad \delta_4>0,
\]
并且对一切 \(4\le q\le 23\) 都有 \(\delta_q>0\)。特别地，
\[
\delta_2=-0.594584\ldots,\qquad
\delta_3=-0.297226\ldots,\qquad
\delta_4=0.002502\ldots.
\]
因此
\[
\mathcal S_q<\infty \qquad (q=2,3),
\]
而从 \(q=4\) 起，对每个 \(4\le q\le 23\) 都有
\[
\sum_{n\le N}\frac{|E_q(n)|^2}{r_q^n}\asymp e^{\delta_q N}.
\]
故四阶碰撞核正是二次残差从平方根稳定相跃迁到边界主导相的最小阶。
\end{theorem}

\begin{proof}
表 \ref{tab:fold_collision_kernel_rh_scan_q2_17} 与表 \ref{tab:fold_collision_kernel_rh_scan_q18_23} 的读数已由注 \ref{rem:pom-collision-kernel-rh-phase} 总结为：在已核验范围 \(q\le 23\) 内，\(\delta_q<0\) 仅发生于 \(q=2,3\)，并且自 \(q=4\) 起严格越线且随后单调上升。上列三个数值即来自同两张表。

再由定理 \ref{thm:pom-deltaq-mean-square-rh-criterion}，
\[
\delta_q<0 \iff \mathcal S_q<\infty,
\qquad
\delta_q>0 \iff \sum_{n\le N}\frac{|E_q(n)|^2}{r_q^n}\asymp e^{\delta_q N}.
\]
代入符号判据便得到结论。证毕。
\end{proof}

\begin{corollary}[中心化二次观测的指数样本壁与已核验窗中的无再入性]\label{cor:conclusion-kernel-rh-sample-wall-no-reentry}
设 \(X_{m,q}\) 为任意满足
\[
\EE[X_{m,q}]=c\,r_q^{m/2},
\qquad
\Var(X_{m,q})=\Theta(\Lambda_q^{2m})
\]
的可采样中心化二次观测，并令 \(\widehat X_{m,q,N}\) 为 \(N\) 次独立样本均值。则相对均方误差满足
\[
\frac{\EE[(\widehat X_{m,q,N}-\EE X_{m,q})^2]}{\EE[X_{m,q}]^2}
\asymp
\frac{1}{N}e^{m\delta_q}.
\]
因此，要使该误差在 \(m\to\infty\) 下趋于 \(0\)，必要且常数意义下充分的样本量标度为
\[
N_q(m)\asymp e^{m\delta_q}.
\]
在已核验窗口 \(4\le q\le 23\) 上，\(\delta_q\) 严格为正且单调上升；特别地，
\[
\delta_4\approx 0.002502,\qquad
\delta_{23}\approx 5.220673.
\]
故样本复杂度在 \(q=4\) 处进入近临界指数壁，并在 \(4\le q\le 23\) 内不再返回平方根相。
\end{corollary}

\begin{proof}
推论 \ref{cor:pom-deltaq-sample-complexity-threshold} 已给出相对均方误差的精确指数尺度
\[
\frac{\Var(X_{m,q})}{N\,\EE[X_{m,q}]^2}\asymp \frac{1}{N}e^{m\delta_q},
\]
并由此推出 \(N_q(m)\asymp e^{m\delta_q}\)。
再由注 \ref{rem:pom-collision-kernel-rh-phase} 与表 \ref{tab:fold_collision_kernel_rh_scan_q18_23}，\(\delta_q\) 在 \(q=4,\dots,23\) 上严格为正并随 \(q\) 上升，故指数壁一经触发便无再入。证毕。
\end{proof}

\leanverified{paper\_conclusion\_kernel\_rh\_codimension\_identity}
\begin{theorem}[越线指数的谱几何余维恒等式]\label{thm:conclusion-kernel-rh-codimension-identity}
定义 bulk 与 bdry 谱维数量纲
\[
d_{\mathrm{bulk}}(q):=\frac{\log r_q}{\log\varphi},
\qquad
d_{\mathrm{bdry}}(q):=\frac{\log\Lambda_q}{\log\varphi}.
\]
则恒有
\[
\frac{\delta_q}{\log\varphi}=2d_{\mathrm{bdry}}(q)-d_{\mathrm{bulk}}(q).
\]
若再定义相关余维
\[
\operatorname{codim}_{\mathrm{corr}}(q):=-\frac{\delta_q}{\log\varphi},
\]
则
\[
\delta_q=0
\iff
d_{\mathrm{bdry}}(q)=\frac12 d_{\mathrm{bulk}}(q)
\iff
\beta_q^{\mathrm{RH}}=\frac12.
\]
因此，kernel-RH 越线恰等价于边界谱维数越过 bulk 半维门槛；它不是模糊的谱隙劣化，而是一个精确的谱几何事件。
\end{theorem}

\begin{proof}
命题 \ref{prop:pom-deltaq-spectral-dimension-identity} 已给出
\[
\frac{\delta_q}{\log\varphi}=2d_{\mathrm{bdry}}(q)-d_{\mathrm{bulk}}(q)
\]
与
\(
\operatorname{codim}_{\mathrm{corr}}(q)=-\delta_q/\log\varphi
\)。
推论 \ref{cor:pom-deltaq-half-threshold-criterion} 又给出
\[
\delta_q=0
\iff
d_{\mathrm{bdry}}(q)=\tfrac12 d_{\mathrm{bulk}}(q)
\iff
\beta_q^{\mathrm{RH}}=\tfrac12.
\]
两式合并即得结论。证毕。
\end{proof}

\begin{theorem}[两步饱和下的一步黄金残差与复缺口的 \texorpdfstring{$2/q$}{2/q} 律]\label{thm:conclusion-kernel-rh-golden-residual-two-step-saturation}
\leanverified{paper\_conclusion\_kernel\_rh\_golden\_residual\_two\_step\_saturation}
设
\[
R_q:=\frac{r_q}{r_{q-2}},
\qquad
\eta_q:=\frac{|\mu_q|}{r_{q-2}},
\qquad
g_q^{(c)}:=1-\frac{\log|\mu_q|}{\log r_q}.
\]
若满足两条饱和假设
\[
\lim_{q\to\infty}r_q^{1/q}=\sqrt{\varphi},
\qquad
\eta_q\longrightarrow 1,
\]
则有严格极限
\[
\frac{|\mu_q|}{r_q}\longrightarrow \varphi^{-1},
\qquad
q\,g_q^{(c)}\longrightarrow 2.
\]
因此，二步尺度上已饱和的复共振模态，在单步 Perron 归一化之后仍留下一个不可消去的黄金残差 \(\varphi^{-1}\)；而其相对指数缺口则精确按 \(2/q\) 退零。
\end{theorem}

\begin{proof}
由推论 \ref{cor:pom-aitken-golden-limit}，在上述两条假设下，
\[
\frac{|\mu_q|}{r_q}
=
\delta_{\mathrm{ait}}(q)
=
\frac{\eta_q}{R_q}
\longrightarrow \varphi^{-1}.
\]
另一方面，推论 \ref{cor:pom-resonance-complex-gap-2overq} 已证明
\[
q\,g_q^{(c)}\longrightarrow 2.
\]
两式合并即得结论。证毕。
\end{proof}
